% Abstract page styled after the NBI project-report convention.
\begin{center}
    {\Large \textcolor{BrickRed}{UNIVERSITY OF COPENHAGEN}\par}
    \vspace{0.9cm}

    {\Huge\itshape Abstract\par}
    \vspace{0.9cm}

    {\large \textcolor{BrickRed}{Faculty of Science}\par}
    {\large \textcolor{BrickRed}{The Niels Bohr Institute}\par}
    \vspace{0.6cm}

    {\large Master Thesis Preparation Project\par}
    \vspace{0.6cm}

    {\bfseries Narrowing the Gap Between Simulation and Real Data in IceCube\par}
    \vspace{0.3cm}

    by Holger Klevang Christiansen\par
\end{center}
\vspace{0.9cm}

Essentially every IceCube machine-learning result rests on the assumption that
simulated and real detector data look sufficiently alike, since the models are
trained on simulation before being applied to data. This report tests that
assumption on a high-statistics sample of atmospheric muons, comparing roughly
two million Monte Carlo events against two million classified muons from the
2021 burnsample. Before the comparison, the report lays the groundwork, the
IceCube detector and its pulse-level data together with the machine-learning
tools the analysis relies on. A pulse-level transformer trained to
separate the two samples serves as the benchmark, and it initially does so
almost perfectly, with test-set AUCs of $0.9882$ for stopped and $0.9960$ for
through-going muons. Four corrections are then applied in sequence, an angular
reweighting of the simulation, a pulse-merging step, a data-driven re-labelling
of HLC pulses, and the removal of low-information events flagged by a learned
von Mises--Fisher uncertainty. Together they lower the AUCs to $0.9218$ and
$0.9848$, with the HLC re-labelling providing the largest single step overall, and they
strip away much of the classifier's per-event confidence. The two samples
nevertheless remain clearly distinguishable. The gap between IceCube simulation
and real data has been narrowed and several of its carriers identified, but it
has not been closed.
